\section{Pregunta N$^{\circ}$12\qquad Carlos Daniel Malvaceda Canales}

\begin{frame}
	\begin{enumerate}\setcounter{enumi}{11}
		\item

		      En una heladería, por $1$ helado, $2$ zumos y $4$ batidos
		      nos cobraron S/$35$.
		      Otro día, por $4$ helados, $4$ zumos y $1$ batido nos
		      cobraron S/$34$.
		      Un tercer día por $2$ helados, $3$ zumos y $4$ batidos
		      S/$42$.
		      Determine el precio de cada uno según el siguiente requerimiento.

		      \begin{enumerate}[a)]
			      \item

			            Modele el sistema.

			      \item
			            Resuelve el sistema usando el método de Parlet y Reid.
		      \end{enumerate}

	\end{enumerate}

	\begin{solution}
		 x : precio del helado , y : precio del zumo , z : precio del batido 
        \begin{align*}
            1x+2y+4z=35 \\
            4x+4y+1z=34 \\ 
            2x+3y+4z=42 \\
        \end{align*}
         \begin{align*}
            \begin{pmatrix}
            1 & 2 & 4 \\
            4 & 4 & 1 \\
            2 & 3 & 4 
            \end{pmatrix} .
            \begin{pmatrix}
                x \\ y \\ z
            \end{pmatrix} = 
            \begin{pmatrix}
                35 \\ 34 \\ 42
            \end{pmatrix}
         \end{align*}
         La matriz no es simétrica , por lo que no se puede aplicar la descomposición Parleit Reidd.
         Realizamos el intercambio de filas para que la matriz sea simetrica.
        
	\end{solution}
\end{frame}

\begin{frame}
	\begin{solution}
		\begin{align*}
            \begin{pmatrix}
            1 & 2 & 4 \\
            2 & 3 & 4 \\
            4 & 4 & 1 
            \end{pmatrix} .
            \begin{pmatrix}
                x \\ y \\ z
            \end{pmatrix} = 
            \begin{pmatrix}
                35 \\ 42 \\ 34
            \end{pmatrix}
         \end{align*}
        
         La descomposición Parlett Reidd tiene que tener la forma $ PAP^{T} = LTL^{T} $ , lo cual nos llevara $n-2$ pasos. \\ 
         La matriz es símetrica e indefinida , por que sus autovalores son negativos y positivos. Ahora realizaremos las iteraciones :

        k = 1 : 
    \begin{equation*}
        A = 
        \begin{bmatrix}
        1 & 2 & 4 \\
            2 & 3 & 4 \\
            4 & 4 & 1 
        
        \end{bmatrix}
    \end{equation*}
    Buscamos el elemento del vector \( [a_{21},a_{31}] \)
    de mayor valor absoluto y se determina una permutación, que representaremos por \(\tilde{P}_{2}\) , tal que :
    \[ \tilde{P}_{1} \begin{bmatrix}
        a_{21} \\ a_{31}
        \end{bmatrix} = \begin{bmatrix}
        \hat{a}_{2} \\ \hat{a}_{3}
        \end{bmatrix} , donde |\hat{a} |_{2} = max [|\hat{a}_{21}| , |\hat{a}_{31}| ] \]
        
	\end{solution}
\end{frame}

\begin{frame}{}
    \begin{solution}
    Donde : \\
        \begin{center}
            \alpha_1 = \begin{bmatrix}
            0 \\ 0 \\ \hat{a}_{3} / \hat{a}_{2}
            \end{bmatrix} , P_1 = [e_1,e_3,e_2] \\
     M_1 = I_3 - \alpha_1 e_2 = I_3 - \begin{bmatrix}
            0 \\ 0 \\ 2/4
            \end{bmatrix} \begin{bmatrix}
            0 & 1 & 0 
            \end{bmatrix} \\ 
        P_1 = \begin{pmatrix}
            1 & 0 & 0 \\ 0 & 0 & 1 \\ 0 & 1 & 0

        \end{pmatrix}   
            
        A^{(1)} = M_1 P_1 A P_1^{T} M_1^{T} = \begin{bmatrix}
        1 & 4 & 0 \\ 
        4 & 1 & 7/2 \\
        0 & 7/2 & -3/4
        \end{bmatrix}
       
        \end{center}

         Ya tenemos la matriz tridiagonal , se cumplieron los n-2 pasos.
       
    \end{solution}
    
\end{frame}

\begin{frame}{Continuación}
        \begin{align*}
            P = P_1 = \begin{pmatrix}
                1 & 0 & 0 \\ 0 & 0 & 1 \\ 0 & 1 & 0
            \end{pmatrix} \\
            T =  M_1 A P_1 ^ {T} M_1 ^ {T}  = \begin{pmatrix}
                1 & 4 & 0 \\ 4 & 1 & 3.5 \\ 0 & 3.5 & -0.75
            \end{pmatrix}
            \\
            L = (M_1 P_1 P^T ) ^{-1}  = \begin{pmatrix}
                1 & 0 & 0 \\ 0 & 1 & 0 \\ 0 & 1/2 & 1
            \end{pmatrix}
            \\
            \begin{bmatrix}
                x & y & z
            \end{bmatrix} = 
            \begin{bmatrix}
                3 & 4 & 6
            \end{bmatrix} 
            \\ 
        \end{align*}
        
        

\end{frame}

\begin{frame}{}
    Mediante una factorización , como esta la resolución $Ax=b$ constaria de las siguientes etapas.
    \begin{enumerate}
        \item L $z$ = P. b
        \item   T W = z (Se resuelve por Thomas)
        \item L$^T$ y = W
        \item x = P$^{T}$y
        
    \end{enumerate}
    
    \begin{align*}
        z = \begin{bmatrix}
            35 & -106 & 0.2666
        \end{bmatrix} \\
        w = \begin{bmatrix}
            3 & 8 & 4
        \end{bmatrix} \\ 
        y = \begin{bmatrix}
            3 & 6 & 4
        \end{bmatrix}
    \end{align*}
    Estos valores estan determinados por el código python adjunto , el cual mediante la resolución de matrices y sustitución se llega a la resolución de [x,y,z]

    Donde los precios son :
\begin{align*}
    x :  helado = 3 \\ y :  zumo  = 4 \\ z : batido = 6
\end{align*}
 
\end{frame}